\documentclass[12pt,a4paper]{article}

% Cấu hình tiếng Việt
\usepackage[utf8]{inputenc}
\usepackage[T5]{fontenc}
\usepackage[vietnamese]{babel}

% Cấu hình căn lề và định dạng
\usepackage{geometry}
\geometry{a4paper, left=2cm, right=2cm, top=2cm, bottom=2cm}
\usepackage{hyperref}
\usepackage{parskip}

% Cấu hình bảng
\usepackage{tabularx}
\usepackage{array}

% Cấu hình hiển thị code
\usepackage{listings}
\usepackage{xcolor}

\definecolor{codegreen}{rgb}{0,0.6,0}
\definecolor{codegray}{rgb}{0.5,0.5,0.5}
\definecolor{codepurple}{rgb}{0.58,0,0.82}
\definecolor{backcolour}{rgb}{0.95,0.95,0.92}

\lstset{
    backgroundcolor=\color{backcolour},   
    commentstyle=\color{codegreen},
    keywordstyle=\color{magenta},
    numberstyle=\tiny\color{codegray},
    stringstyle=\color{codepurple},
    basicstyle=\ttfamily\footnotesize,
    breakatwhitespace=false,         
    breaklines=true,                 
    captionpos=b,                    
    keepspaces=true,                 
    numbers=left,                    
    numbersep=5pt,                  
    showspaces=false,                
    showstringspaces=false,
    showtabs=false,                  
    tabsize=2
}

\title{\textbf{Buổi 6: Authentication và Authorization}}
\author{}
\date{}

\begin{document}

\maketitle

Trong bối cảnh \textbf{Kiến trúc hướng dịch vụ (Service-Oriented Architecture - SOA)} hoặc Microservices, việc quản lý định danh trở nên phức tạp hơn rất nhiều so với kiến trúc Monolithic (nguyên khối). Thay vì một server duy nhất quản lý session của người dùng, chúng ta có hàng chục, hàng trăm service độc lập cần biết \textit{``Người đang gọi request này là ai?''} và \textit{``Họ có quyền làm việc này không?''}.

Dưới đây là giải thích chi tiết cho các mục tiêu kiến thức và kỹ năng của buổi học này.

\vspace{0.5cm}
\hrule
\vspace{0.5cm}

\section{So sánh JWT vs OAuth 2.0}

Rất nhiều người nhầm lẫn giữa JWT và OAuth 2.0 và coi chúng là hai công nghệ cạnh tranh nhau. Thực tế, chúng giải quyết hai bài toán khác nhau và thường được \textbf{sử dụng cùng nhau}.

\begin{itemize}
    \item \textbf{OAuth 2.0} là một \textbf{Giao thức phân quyền (Authorization Protocol / Framework)}. Nó định nghĩa các luồng (flows/grants) để một ứng dụng (Client) có thể xin phép người dùng để truy cập vào tài nguyên của họ trên một máy chủ khác (Resource Server) mà không cần giữ mật khẩu của người dùng.
    \item \textbf{JWT (JSON Web Token)} là một \textbf{Định dạng dữ liệu (Data Format)}. Nó là một chuỗi ký tự được mã hóa và ký điện tử (signed) chứa các thông tin (claims) về người dùng hoặc hệ thống.
\end{itemize}

\subsection*{Bảng so sánh sự khác biệt}

\begin{table}[h]
\centering
\renewcommand{\arraystretch}{1.5}
\begin{tabular}{|p{3cm}|p{6cm}|p{6cm}|}
\hline
\textbf{Tiêu chí} & \textbf{JWT (JSON Web Token)} & \textbf{OAuth 2.0} \\ 
\hline
\textbf{Bản chất} & Là một định dạng token (phương tiện mang thông tin). & Là một bộ khung/giao thức (quy trình hoạt động). \\ 
\hline
\textbf{Mục đích chính} & Xác thực tĩnh (Stateless Authentication) và truyền tải thông tin an toàn. & Ủy quyền truy cập (Delegated Authorization) an toàn. \\ 
\hline
\textbf{Trạng thái} & Stateless (Không lưu trạng thái trên server). Mọi thông tin nằm trong token. & Stateful hoặc Stateless (tùy thuộc vào cách IdP quản lý token). \\ 
\hline
\textbf{Mối quan hệ} & JWT thường được dùng làm định dạng cho Access Token bên trong luồng OAuth 2.0. & OAuth 2.0 sinh ra token, và token đó có thể là JWT. \\ 
\hline
\end{tabular}
\end{table}

\vspace{0.5cm}
\hrule
\vspace{0.5cm}

\section{Hiểu các khái niệm cốt lõi}

Trong kiến trúc SOA, khi một Client gọi qua API Gateway để vào các dịch vụ nội bộ (ví dụ: Order Service, Payment Service), các khái niệm sau sẽ đóng vai trò kiểm soát bảo mật.

\subsection{Bearer Token}
\textbf{Bearer} có nghĩa là ``người mang/người sở hữu''. Bearer token là một loại token mà bất kỳ ai nắm giữ nó đều có thể sử dụng nó để truy cập tài nguyên.
\begin{itemize}
    \item \textbf{Cách sử dụng:} Gắn vào HTTP Header của mỗi request: \texttt{Authorization: Bearer <token\_string>}
    \item \textbf{Đặc điểm:} Server không quan tâm ai là người gửi request, chỉ cần request có chứa Bearer Token hợp lệ, server sẽ cấp quyền truy cập. Do đó, việc bảo mật token này là tối quan trọng.
\end{itemize}

\subsection{Refresh Token}
Access token (thường là JWT) cần có tuổi thọ rất ngắn (ví dụ: 15 phút) để giảm thiểu rủi ro nếu bị lộ. Khi Access Token hết hạn, hệ thống sử dụng \textbf{Refresh Token} để xin cấp lại một Access Token mới mà không bắt người dùng phải đăng nhập lại.
\begin{itemize}
    \item \textbf{Đặc điểm:} Tuổi thọ dài (vài ngày, vài tháng). Thường được lưu trữ an toàn trong cơ sở dữ liệu của server (Stateful) để có thể thu hồi (revoke) bất cứ lúc nào.
\end{itemize}

\subsection{Scopes (Phạm vi)}
Scope là khái niệm thuộc về \textbf{OAuth 2.0}. Nó định nghĩa quyền hạn của \textbf{Ứng dụng (Client Application)}.
\begin{itemize}
    \item \textbf{Ví dụ:} Khi bạn dùng tài khoản Google để đăng nhập vào một trang web tìm việc, trang web đó (Client) sẽ xin cấp quyền (Scope) là \texttt{read:email} và \texttt{read:profile}. Nó không có quyền đọc Google Drive của bạn vì không nằm trong scope.
\end{itemize}

\subsection{Roles (Vai trò)}
Role là khái niệm thuộc về \textbf{RBAC (Role-Based Access Control)}. Nó định nghĩa quyền hạn của \textbf{Người dùng (User)}.
\begin{itemize}
    \item \textbf{Ví dụ:} Trong hệ thống nội bộ, người dùng Nguyễn Văn A có role là \texttt{Admin}, người dùng B có role là \texttt{Customer}. Dựa vào Role này, hệ thống SOA sẽ quyết định (Authorization) họ được gọi tới API nào.
\end{itemize}

\vspace{0.5cm}
\hrule
\vspace{0.5cm}

\section{Kỹ năng: Triển khai JWT Authentication cho một API đơn giản}

Trong SOA, luồng xác thực JWT cơ bản diễn ra như sau:
\begin{enumerate}
    \item Client gửi thông tin đăng nhập (User/Pass) tới \textbf{Authentication Service}.
    \item Service kiểm tra đúng sai, tạo ra một chuỗi JWT và trả về cho Client.
    \item Client lưu JWT. Lần sau gọi API, Client đính kèm JWT vào header \texttt{Authorization: Bearer <JWT>}.
    \item API Gateway sẽ verify (xác minh) chữ ký của JWT. Nếu hợp lệ, cho phép request đi tiếp.
\end{enumerate}

Dưới đây là mô phỏng logic bằng \textbf{Node.js/Express} (chỉ tập trung vào phần verify):

\begin{lstlisting}[language=JavaScript]
const jwt = require('jsonwebtoken');
const SECRET_KEY = "chuoi_ky_mat_cua_he_thong";

// Middleware xac thuc JWT
function authenticateToken(req, res, next) {
    // 1. Lay token tu header
    const authHeader = req.headers['authorization'];
    const token = authHeader && authHeader.split(' ')[1]; // Format: "Bearer <token>"

    // 2. Neu khong co token -> Tu choi
    if (!token) return res.status(401).json({ message: "Khong tim thay token" });

    // 3. Xac minh chu ky va thoi han cua token
    jwt.verify(token, SECRET_KEY, (err, user) => {
        if (err) return res.status(403).json({ message: "Token khong hop le hoac da het han" });
        
        // 4. Luu thong tin user (claims) vao request de cac ham sau su dung
        req.user = user; 
        next();
    });
}

// API can bao ve
app.get('/api/v1/orders', authenticateToken, (req, res) => {
    res.json({ message: `Lay danh sach don hang cho user ID: ${req.user.id}` });
});
\end{lstlisting}

\vspace{0.5cm}
\hrule
\vspace{0.5cm}

\section{Kỹ năng: Phát hiện và khắc phục rủi ro bảo mật}

\subsection{Rủi ro 1: Token Leakage (Rò rỉ Token)}
\begin{itemize}
    \item \textbf{Bản chất:} Kẻ tấn công lấy cắp được JWT của người dùng hợp lệ.
    \item \textbf{Nguyên nhân phổ biến:} Lưu JWT trong \texttt{localStorage} hoặc \texttt{sessionStorage} trên trình duyệt, dễ bị đánh cắp qua lỗ hổng XSS; hoặc gửi token qua kết nối HTTP không mã hóa.
    \item \textbf{Cách khắc phục:}
    \begin{itemize}
        \item \textbf{Sử dụng HttpOnly Cookies:} Đẩy JWT vào Cookie và set cờ \texttt{HttpOnly} + \texttt{Secure}. Trình duyệt sẽ tự động đính kèm cookie này vào các request mà JavaScript ở frontend không thể đọc được.
        \item \textbf{Bắt buộc dùng HTTPS:} Mã hóa đường truyền bằng TLS/SSL.
    \end{itemize}
\end{itemize}

\subsection{Rủi ro 2: Replay Attack (Tấn công phát lại)}
\begin{itemize}
    \item \textbf{Bản chất:} Kẻ tấn công chặn bắt được một request chứa Bearer Token hợp lệ, sau đó gửi đi gửi lại request đó nhiều lần để thực hiện các hành vi trục lợi.
    \item \textbf{Cách khắc phục:}
    \begin{itemize}
        \item \textbf{Short-lived Access Token:} Cài đặt thời gian sống (Expiration - \texttt{exp}) của JWT cực ngắn (5 - 15 phút).
        \item \textbf{Sử dụng JWT ID (\texttt{jti} claim):} Mỗi JWT sinh ra có một ID duy nhất. Hệ thống sẽ lưu lại danh sách các \texttt{jti} đã sử dụng để từ chối các request lặp lại.
        \item \textbf{Thêm Context vào Token:} Gắn thêm địa chỉ IP hoặc thông tin thiết bị vào payload của JWT để đối chiếu khi verify.
    \end{itemize}
\end{itemize}

\end{document}